\section{问题重述}

\subsection{问题背景}

无创产前检测（Non-invasive Prenatal Test，NIPT）是一种通过采集孕妇血液并分析其中胎儿游离DNA片段，对胎儿染色体异常风险进行评估的产前检测技术。与传统侵入性产前诊断方式相比，NIPT具有安全性较高、检测过程方便等特点，在唐氏综合征、爱德华氏综合征和帕陶氏综合征等常见染色体异常的早期筛查中具有重要作用。

NIPT检测结果的可靠性受到孕周、孕妇身体质量指数以及测序质量等多种因素影响。对于男胎，当Y染色体浓度达到一定水平时，检测结果才具有较高的可信度。若检测时间过早，胎儿游离DNA浓度可能不足，从而增加检测失败或结果不稳定的可能；若检测时间过晚，又可能缩短发现异常后的临床处理时间。因此，如何在检测可靠性与早期发现风险之间取得平衡，是NIPT时点选择中的重要问题。

此外，不同孕妇在年龄、身高、体重、BMI和妊娠状况等方面存在个体差异，采用统一的检测时间难以适用于所有人群。对于不携带Y染色体的女胎，还需要结合染色体Z值、GC含量和读段质量等信息建立异常判定方法。基于上述背景，本文利用给定的男胎与女胎NIPT检测数据，研究男胎检测时点的合理选择，并建立女胎染色体异常判定模型。

\subsection{问题要求}

针对附件所给出的男胎与女胎NIPT检测数据，本文需要完成以下四项任务：

\textbf{问题一：}分析男胎Y染色体浓度与孕妇检测孕周、BMI及其他相关指标之间的关系，建立相应的关系模型，并检验模型及相关变量的显著性。

\textbf{问题二：}在BMI是影响男胎Y染色体浓度最早达标时间主要因素的条件下，对男胎孕妇进行合理的BMI分组，给出各组BMI区间和最佳NIPT检测时点，使潜在风险尽量降低，并分析检测误差的影响。

\textbf{问题三：}在问题二基础上，进一步考虑年龄、身高、体重等个体因素，以及检测误差和Y染色体浓度达标比例的影响，重新确定合理的BMI分组和各组最佳NIPT检测时点。

\textbf{问题四：}针对女胎不携带Y染色体的特点，以附件AB列中的13、18、21号染色体非整倍体结果作为异常判定依据，综合考虑染色体Z值、GC含量、测序读段指标、相关比例、X染色体浓度和BMI等因素，建立女胎异常判定方法。

\section{问题分析}

\subsection{问题一分析}

问题一研究男胎Y染色体浓度与检测孕周、BMI及其他指标之间的数量关系，是后续确定最佳检测时点的基础。首先统一孕周表示方式并进行描述性统计和可视化，再建立Y染色体浓度与各影响因素之间的关系模型，通过参数显著性、拟合优度和残差分析评价模型。由于同一孕妇可能存在多次检测记录，还需处理个体内相关性，预期得到孕周、BMI等因素的影响方向、影响程度及统计显著性。

\subsection{问题二分析}

问题二承接问题一，将研究目标由“解释Y染色体浓度的变化”转化为“确定不同BMI人群的最佳NIPT检测时点”。利用问题一中孕周、BMI与Y染色体浓度之间的关系，将Y染色体浓度达到或超过$4\%$定义为达标，估计不同BMI和孕周下的达标概率。在此基础上对BMI进行数据驱动分组，并在检测可靠性与延迟发现风险之间进行权衡，为每组选择尽可能早且可靠的检测时点。最后通过重复测序、Bootstrap和误差扰动分析结论的稳定性。

\subsection{问题三分析}

问题三是对问题二的扩展。问题二主要依据BMI进行群体分组，但实际达标时间还可能受到年龄、身高、体重、妊娠方式和测序质量的影响。因此需在保留BMI分组这一实际应用形式的同时，建立多因素达标概率模型，将更多个体特征纳入组内风险计算，重新确定各组最佳检测时点，并检验误差条件下分组边界和推荐时点的稳定性。

\subsection{问题四分析}

问题四的研究对象由男胎转为女胎。本文将附件AB列为空定义为无异常，出现T13、T18、T21或其组合定义为异常；以各染色体Z值、X染色体浓度、GC含量、测序读段指标、BMI和孕周等作为候选特征建立分类模型。训练集和测试集按孕妇代码划分，避免重复记录造成信息泄漏，并使用召回率、特异度、F1值和ROC--AUC等指标评价女胎异常判定效果。
